% Copyright (C) 2026 臧炫懿
% 
% This program is free software: you can redistribute it and/or modify
% it under the terms of the GNU General Public License as published by
% the Free Software Foundation, either version 3 of the License, or
% (at your option) any later version.
%
% This program is distributed in the hope that it will be useful,
% but WITHOUT ANY WARRANTY; without even the implied warranty of
% MERCHANTABILITY or FITNESS FOR A PARTICULAR PURPOSE. See the
% GNU General Public License for more details.
%
% You should have received a copy of the GNU General Public License
% along with this program. If not, see <https://www.gnu.org/licenses/>.

\usepackage[fontset=none]{ctex}
\usepackage{fontspec}
\usepackage{minted}
\usepackage{graphicx}
\usepackage{multicol}
\usepackage{hyperref}

% \setlength{\parindent}{2em}
\setlength{\parskip}{0.5em}

\setminted{
    fontsize=\small,           % 字体大小
    baselinestretch=1.0,       % 行间距
    linenos=true,              % 行号
    frame=lines,               % 边框样式
    framesep=2mm,              % 边框间距
    rulecolor=\color{black},   % 边框颜色
    numbersep=5pt,             % 行号间距
    breaklines=true,           % 自动换行
    breakanywhere=false,       % 不在任意位置断行
    showspaces=false,          % 不显示空格
    showtabs=false,            % 不显示制表符
    tabsize=4,                 % 制表符大小
    % url=false,
}


\setmainfont{Source Serif Pro}
\setsansfont{Source Sans Pro}[Scale=MatchLowercase]
\setmonofont{Source Code Pro}[Scale=MatchLowercase]
\setCJKmainfont[ItalicFont = * ]{Source Han Serif CN}
\setCJKsansfont[ItalicFont = * ]{Source Han Sans CN}
\setCJKmonofont[ItalicFont = * ]{Source Han Sans CN}

\usetheme{Berlin}
\usefonttheme{serif}

\title{《AI基础》课程编程指南}
\author{臧炫懿}
\date{\today}